%% toc: P3 31 31.1,31.2,31.3,31.4
\chapter{Indiscernibles, Ehrenfeucht--Mostowski Models, and Cardinal Transfer}
\label{sec:p03-w03-p03-028-chapter}

\Cref{sec:p03-indiscernibles} defines order indiscernibles and states the
Ehrenfeucht--Mostowski theorems, and
\cref{sec:p03-morley-ramsey-models} proves the indiscernible and term-hull
construction in the form Morley gave it.  This chapter supplies the
combinatorial input that both rely on, Ramsey's partition theorems, proves
them in Ramsey's own formulation, and fixes exactly where a partition
theorem enters the construction: with which arity, with how many colours,
and in which step.  It then records the whole compactness, Ramsey and
Skolem-hull route step by step, with the source and locator for each step,
and separates what the route proves from the stronger spectrum statement of
\cref{thm:p03-unstable-spectrum}.

\emph{Purpose.}  To make the existence of order indiscernibles a proved
consequence of first-order compactness and one named partition theorem, and
to state the boundary of each source used.

\emph{Inputs.}  First-order compactness (\cref{thm:fol-compactness});
the axiom of choice, which Ramsey's infinite theorem assumes; order
indiscernibles (\cref{def:p03-order-indiscernibles}); Skolem expansions and
term hulls (\cref{def:p03-morley-skolem-hull}).

\emph{Conclusions.}  Ramsey's Theorems A, B and C
(\cref{thm:p03-w03-p03-028-ramsey-a,thm:p03-w03-p03-028-ramsey-c,thm:p03-w03-p03-028-ramsey-b});
the partition step of the indiscernible construction
(\cref{thm:p03-w03-p03-028-ramsey-step}) and the finite satisfiability of
the indiscernibility diagram (\cref{thm:p03-w03-p03-028-finite-sat});
coherence of the blueprint read off an indiscernible sequence
(\cref{thm:p03-w03-p03-028-coherent}) and its realization over any linear
order (\cref{thm:p03-w03-p03-028-realize}), with the elementary
embeddings induced by order embeddings
(\cref{thm:p03-w03-p03-028-elementary}); the Ehrenfeucht--Mostowski
automorphism theorems (\cref{thm:p03-w03-p03-028-em-57,thm:p03-w03-p03-028-em-56})
and their limit (\cref{thm:p03-w03-p03-028-no-finite-order}).  Every result
of this chapter is proved in Annex~\ref{sec:annex-p03-w03-p03-028}.

\section{Ramsey theorems and indiscernible sequences}
\label{sec:p03-w03-p03-028-ramsey}

Ramsey proved his partition theorems as a preliminary to a decision
procedure for a class of first-order formulas.  The theorem his logical
argument needs concerns finite classes (Theorem~B); he first proves an
infinite analogue (Theorem~A) because its proof is simpler and shows the
method \cite[p.~264]{Ramsey1930}.  The indiscernible construction uses the
infinite theorem; the finite one is proved here as well, because it gives
an alternative route through the same step
(\cref{def:p03-w03-p03-028-finite-route}).

\begin{hbconvention}{def:p03-w03-p03-028-ramsey-notation}{Ramsey's notation}
The notation follows \cite[pp.~264, 266--267]{Ramsey1930}.  An
\emph{\(r\)-combination} of a class \(\Gamma\) is a subclass of \(\Gamma\)
with exactly \(r\) members.  To \emph{divide} the \(r\)-combinations of
\(\Gamma\) into \(\mu\) mutually exclusive classes \(C_i\)
(\(i=1,\ldots,\mu\)) is to assign each \(r\)-combination to exactly one
\(C_i\); some \(C_i\) may be empty.  In current terminology this is a
colouring with \(\mu\) colours, and a subclass all of whose
\(r\)-combinations lie in one \(C_i\) is \emph{homogeneous}.  For finite
classes, a suffix records the number of members: \(\Gamma_m\) is a class
with \(m\) members and \(\Delta_n\) one with \(n\) members.  The letters
\(r,n,k,m,\mu,\nu\) denote finite cardinals, and \(r,n,\mu\) are positive.
Suffixes on \(C\) only distinguish classes of combinations and do not count
members.  When \(\Gamma=\omega\) the statement ``every division of the
\(r\)-combinations into \(\mu\) classes has an infinite homogeneous
subclass'' is the partition relation
\(\omega\longrightarrow(\omega)^r_\mu\) of \cref{def:p06-partition-arrow}.
\end{hbconvention}

\begin{hbtheorem}{thm:p03-w03-p03-028-ramsey-a}{proved=annex:p03-w03-p03-028-ramsey-a}{Ramsey's Theorem A}
Let \(\Gamma\) be an infinite class and let \(\mu\) and \(r\) be positive
integers.  Suppose the \(r\)-combinations of the members of \(\Gamma\) are
divided in any manner into \(\mu\) mutually exclusive classes \(C_i\)
(\(i=1,\ldots,\mu\)).  Assuming the axiom of choice, \(\Gamma\) contains an
infinite subclass \(\Delta\) all of whose \(r\)-combinations belong to the
same \(C_i\) \cite[p.~264, proof pp.~264--266]{Ramsey1930}.
\end{hbtheorem}

With \(\Gamma=\omega\) this is \cref{thm:p06-infinite-ramsey}, and the
first sentence of \cref{thm:p03-morley-ramsey-em} is the same theorem in
Morley's form \cite[Theorem~3.1, p.~525]{Morley1965}.  Ramsey's
statement carries the choice hypothesis explicitly; his proof selects
infinitely many elements and infinite subclasses one after another.  The
proof written for \cref{thm:p03-morley-ramsey-em} reaches the same
conclusion by a different argument: it colours each chosen element by the
class of its homogeneous tail and then keeps the elements of one colour.
Ramsey's own argument, reproduced in the annex to this chapter, splits
instead into the two possibilities that a selection favouring the first
class goes on for ever or stops, and treats more than two classes by
merging two of them.

For the finite theorems Ramsey defines explicit bounds by a double
recursion.

\begin{hbdefinition}{def:p03-w03-p03-028-ramsey-bounds}{Ramsey's bound functions}
Define \(f(r,n,k)\) for positive \(r\) and \(n\) and for \(k\geq0\), and
\(g(r,n,k)\) for \(r\geq2\), \(n\geq0\) and \(k\geq0\), as follows
\cite[p.~267]{Ramsey1930}.  For \(r=1\),
\[
  f(1,n,k)=\max(2n-1,\,n+k)\qquad(n\geq1,\ k\geq0).
\]
For \(r\geq2\), assuming \(f(r-1,\cdot,\cdot)\) defined,
\begin{align*}
  f(r,1,k)&=f(r-1,\,k-r+2,\,r-2)+1 &&(k+1\geq r),\\
  g(r,0,k)&=\max(r-1,\,k),\\
  g(r,n,k)&=f\bigl(r,1,g(r,n-1,k)\bigr) &&(n\geq1),\\
  f(r,n,k)&=f\bigl(r,n-1,g(r,n,k)\bigr) &&(n>1).
\end{align*}
For \(n\geq r\) put \cite[p.~269]{Ramsey1930}
\[
  h(r,n,2)=f(r,\,n-r+1,\,r-1),\qquad
  h(r,n,\mu)=h\bigl(r,\,h(r,n,\mu-1),\,2\bigr)\quad(\mu>2).
\]
For \(n<r\), or \(\mu=1\), put \(h(r,n,\mu)=n\).  The last clause is not in
Ramsey's text: there every \(\Delta_n\) is trivially homogeneous, and the
clause makes \(h\) total.
\end{hbdefinition}

\begin{hblemma}{thm:p03-w03-p03-028-bounds-defined}{proved=annex:p03-w03-p03-028-bounds-defined}{The bound functions are well defined}
The recursion of \cref{def:p03-w03-p03-028-ramsey-bounds} defines
\(g(r,n,k)\) for all \(r\geq2\), \(n\geq0\), \(k\geq0\), and
\(f(r,n,k)\) for all \(r,n\geq1\) and \(k\geq0\) with \(n+k\geq r\) when
\(n>1\) and \(k+1\geq r\) when \(n=1\).  Moreover
\(f(r,n,k)\geq n+k\), \(g(r,n,k)\geq\max(r-1,k)+n\), and
\(h(r,n,\mu)\geq n\).
\end{hblemma}

\begin{hbtheorem}{thm:p03-w03-p03-028-ramsey-c}{proved=annex:p03-w03-p03-028-ramsey-c}{Ramsey's Theorem C}
Let \(r,n,k\) satisfy \(n+k\geq r\), and \(k+1\geq r\) when \(n=1\), and
let \(m_0=f(r,n,k)\).  If \(m\geq m_0\) and the \(r\)-combinations of a
class \(\Gamma_m\) are divided into two mutually exclusive classes \(C_1\)
and \(C_2\), then \(\Gamma_m\) contains two mutually exclusive subclasses
\(\Delta_n\) and \(\Delta_k\) such that all the \(r\)-combinations of
members of \(\Delta_n+\Delta_k\) that contain at least one member of
\(\Delta_n\) belong to the same \(C_i\)
\cite[p.~267, proof pp.~267--269]{Ramsey1930}.
\end{hbtheorem}

Ramsey states Theorem~C for all \(n+k\geq r\); the extra condition for
\(n=1\) is the side condition \(k+1\geq r\) of his recursion clause for
\(f(r,1,k)\).  For \(n=1\) and \(k<r-1\) the conclusion holds trivially,
since then no \(r\)-combination of \(\Delta_1+\Delta_k\) exists.

\begin{hbtheorem}{thm:p03-w03-p03-028-ramsey-b}{proved=annex:p03-w03-p03-028-ramsey-b}{Ramsey's Theorem B}
Given any \(r\), \(n\) and \(\mu\), put \(m_0=h(r,n,\mu)\).  If
\(m\geq m_0\) and the \(r\)-combinations of a class \(\Gamma_m\) are
divided in any manner into \(\mu\) mutually exclusive classes \(C_i\)
(\(i=1,\ldots,\mu\)), then \(\Gamma_m\) contains a subclass \(\Delta_n\)
all of whose \(r\)-combinations belong to the same \(C_i\)
\cite[p.~267, proof pp.~267--269]{Ramsey1930}.
\end{hbtheorem}

Ramsey remarks that he cannot determine how large \(m_0\) must be and
expects his values to be far too large \cite[p.~267]{Ramsey1930}.  For
\(r=\mu=2\) his recursion gives \(2^{n(n-1)/2}\), while a direct argument
gives \(n!\) \cite[p.~269]{Ramsey1930}.  None of the constructions below
depends on the size of the bound, only on its existence.

\subsection{Where the partition theorem enters}
\label{sec:p03-w03-p03-028-step}

The indiscernible construction needs one partition step: in an infinite
model, find an infinite set on which finitely many formulas cannot tell
increasing tuples of one fixed length apart.  The ordering used to say
``increasing'' is chosen arbitrarily and need not be definable.

\begin{hblemma}{thm:p03-w03-p03-028-ramsey-step}{proved=annex:p03-w03-p03-028-ramsey-step}{The partition step}
Let \(\mathcal M\) be an \(\Lang\)-structure with infinite domain, let
\(\lhd\) be any linear order of \(\dom\mathcal M\), let \(N\geq1\), and
let \(\psi_1,\ldots,\psi_q\) be \(\Lang\)-formulas whose free variables are
among \(v_0,\ldots,v_{N-1}\).  There is an infinite
\(Y\subseteq\dom\mathcal M\) such that for every \(j\leq q\) and every two
\(\lhd\)-increasing \(N\)-tuples \(\bar a\) and \(\bar b\) from \(Y\),
\[
  \mathcal M\entails\psi_j(\bar a)
  \quad\Longleftrightarrow\quad
  \mathcal M\entails\psi_j(\bar b).
\]
The proof applies Theorem~A once, with arity \(r=N\) and with
\(\mu=2^q\) classes.
\end{hblemma}

\begin{hblemma}{thm:p03-w03-p03-028-finite-sat}{proved=annex:p03-w03-p03-028-finite-sat}{Finite satisfiability of the indiscernibility diagram}
Let \(T\) be an \(\Lang\)-theory with an infinite model and let \((X,<)\)
be a linear order.  Expand \(\Lang\) by new distinct constants \(c_x\)
(\(x\in X\)) and let \(\Sigma_X\) consist of \(T\), the sentences
\(c_x\neq c_y\) for distinct \(x,y\in X\), and the sentences
\[
  \psi(c_{x_1},\ldots,c_{x_n})\leftrightarrow\psi(c_{y_1},\ldots,c_{y_n})
\]
for every \(n\geq1\), every \(\Lang\)-formula \(\psi\) with free variables
among \(v_0,\ldots,v_{n-1}\), and all \(x_1<\cdots<x_n\) and
\(y_1<\cdots<y_n\) in \(X\).  Every finite subset of \(\Sigma_X\) has a
model.  A finite subset with \(q\) sentences of the last kind, whose
lengths are at most \(N\), is satisfied after one application of
\cref{thm:p03-w03-p03-028-ramsey-step} with these \(N\) and \(q\).
\end{hblemma}

By compactness (\cref{thm:fol-compactness}) \(\Sigma_X\) has a model,
and the interpretations of the constants are distinct order indiscernibles
indexed by \(X\) in the sense of \cref{def:p03-order-indiscernibles}.  This
proves \cref{thm:p03-ehrenfeucht-mostowski}; the same argument is the
second sentence of \cref{thm:p03-morley-ramsey-em}, whose annex handles the
different tuple lengths by one application per length instead of by the
padding used here.

The partition theorem used is therefore Theorem~A, at one finite arity
\(N\) and with finitely many classes \(2^q\); neither infinitely many
classes nor infinite arity is needed, and no stronger partition relation
is used.  Ehrenfeucht and Mostowski make the same move in the same
place.  They order the infinite set \(Y\) of a model by an arbitrary
relation that in general has nothing to do with the relations definable in
the theory \cite[p.~61]{EhrenfeuchtMostowski1956}, divide the
\((2n+1)\)-element subsets of \(Y\) into \(2^s\) classes by the truth values
of \(s\) matrices at their increasing enumerations, and apply their
Theorem~5.1 to obtain an infinite homogeneous \(Y_1\)
\cite[pp.~61--62]{EhrenfeuchtMostowski1956}; in the proof of their
Theorem~5.7 the subsets have \(n\) elements and the classes are again
\(2^s\) in number \cite[p.~64]{EhrenfeuchtMostowski1956}.

\begin{hbmethod}{def:p03-w03-p03-028-finite-route}{The finite-partition alternative}
In the proof of \cref{thm:p03-w03-p03-028-finite-sat} only finitely many
homogeneous elements are used.  Theorem~B can therefore replace Theorem~A:
any \(h(N,\,p+N,\,2^q)\) elements of the infinite model, ordered
arbitrarily, contain a homogeneous subclass with \(p+N\) members, where
\(p\) is the number of constants in the finite subset.  This route does not
use the axiom of choice for the partition step.  It is the route that the
phrase ``finite Ramsey'' describes; the proof written in
\cref{sec:annex-p03-w03-p03-028} and the proofs of
Ehrenfeucht--Mostowski and Morley use Theorem~A.
\end{hbmethod}

\begin{hbexample}{def:p03-w03-p03-028-parity-example}{A partition step on a parity graph}
Let \(\Lang=\{E\}\) and let \(\mathcal M\) be the graph on \(\omega\) with
\(i\mathrel E j\) iff \(i+j\) is odd.  Take \(\lhd\) to be the usual order,
\(N=2\), \(q=1\) and \(\psi_1=E(v_0,v_1)\).  The two classes are the pairs
with odd sum (edges) and the pairs with even sum (non-edges).  The even
numbers form an infinite homogeneous set, lying in the non-edge class.  No
homogeneous set with three or more members lies in the edge class: among any
three numbers two have the same parity and form a non-edge.  The class that
Theorem~A delivers is therefore determined by the colouring, not chosen in
advance.  Using the even numbers in the proof of
\cref{thm:p03-w03-p03-028-finite-sat} for \(T=\operatorname{Th}(\mathcal
M)\) gives indiscernibles that are pairwise non-adjacent.
\end{hbexample}

\begin{hbwarning}{def:p03-w03-p03-028-ramsey-limits}{Limits of the partition step}
Theorem~A fails with infinitely many classes: put each \(1\)-combination of
\(\omega\) in its own class; then no homogeneous subclass has two members.
Theorem~A also gives homogeneity only for \(\lhd\)-increasing tuples of one
length.  In \cref{thm:p03-w03-p03-028-finite-sat} the indiscernibility
asserted is order indiscernibility, not total indiscernibility: in a model
of dense linear order the formula \(v_0<v_1\) holds of every increasing
pair of the indiscernibles produced and fails of the reversed pair.
\end{hbwarning}

\section{Skolem expansions and coherent blueprints}
\label{sec:p03-w03-p03-028-blueprints}

The blueprint of \cref{sec:p03-indiscernibles} is a family of quantifier-free
types in a Skolem language.  A sequence of indiscernibles and a blueprint
are different objects: the first lives in one model, the second prescribes
the increasing types for every index order.  The next lemma shows how an
indiscernible sequence in a Skolem expansion yields a coherent blueprint,
and the one after it constructs a sequence realizing a given blueprint over
any linear order.  The Skolem expansion \(T^{\mathrm{sk}}\) is the
conservative universal expansion of \cref{def:p03-morley-skolem-hull},
whose construction is proved in \cref{thm:p03-morley-ramsey-em}.

\begin{hblemma}{thm:p03-w03-p03-028-coherent}{proved=annex:p03-w03-p03-028-coherent}{An indiscernible sequence determines a coherent blueprint}
Let \(\mathcal N\entails T^{\mathrm{sk}}\), let \((I,<)\) be an infinite
linear order, and let \((a_i)_{i\in I}\) be a sequence of distinct elements
of \(\mathcal N\) that is order indiscernible for quantifier-free
\(\Lang^{\mathrm{sk}}\)-formulas.  For \(n\geq1\) let \(p_n\) be the set
of quantifier-free \(\Lang^{\mathrm{sk}}\)-formulas
\(\varphi(v_1,\ldots,v_n)\) with
\(\mathcal N\entails\varphi(a_{i_1},\ldots,a_{i_n})\) for some
\(i_1<\cdots<i_n\) in \(I\).  Then:
\begin{enumerate}
\item \(p_n\) does not depend on the chosen increasing tuple, and it is a
complete quantifier-free \(n\)-type containing \(v_s\neq v_t\) for
\(s\neq t\);
\item the family \(\Phi=(p_n)_{n\geq1}\) is coherent: if
\(\mathcal K\entails T^{\mathrm{sk}}\) and \((b_1,\ldots,b_n)\) realizes
\(p_n\) in \(\mathcal K\), then for all \(1\leq k_1<\cdots<k_m\leq n\) the
tuple \((b_{k_1},\ldots,b_{k_m})\) realizes \(p_m\).
\end{enumerate}
\end{hblemma}

\begin{hblemma}{thm:p03-w03-p03-028-realize}{proved=annex:p03-w03-p03-028-realize}{Realizing a coherent blueprint over any order}
Let \(\Phi=(p_n)_{n\geq1}\) be obtained as in
\cref{thm:p03-w03-p03-028-coherent} and let \((J,<)\) be any linear order.
There are \(\mathcal K\entails T^{\mathrm{sk}}\) and distinct
\((b_j)_{j\in J}\) in \(\mathcal K\) such that every increasing
\(n\)-tuple \(b_{j_1},\ldots,b_{j_n}\) realizes \(p_n\).  The term hull of
\(\{b_j:j\in J\}\) in \(\mathcal K\) is the universe of a model of
\(T^{\mathrm{sk}}\), and its \(\Lang\)-reduct is a model of \(T\).  This
reduct is \(\operatorname{EM}(J,\Phi)\) of \cref{sec:p03-indiscernibles}.
\end{hblemma}

The quantifier-free diagram of the hull is fixed by \(\Phi\) and the order
type of \(J\): a quantifier-free statement about term values is a
quantifier-free statement about the finitely many generators they use, and
these form an increasing tuple.  The next lemma makes this precise and
shows why the Skolem expansion is taken after Morleyization.

\begin{hblemma}{thm:p03-w03-p03-028-elementary}{proved=annex:p03-w03-p03-028-elementary}{Order embeddings induce elementary embeddings}
Let \(T^*\) be the Morleyization of \(T\): the language gets a new relation
symbol \(R_\varphi\) for each \(\Lang\)-formula \(\varphi(\bar x)\), with
the axioms \(\forall\bar x\,(R_\varphi(\bar x)\leftrightarrow\varphi(\bar
x))\).  Let \(T^{*\mathrm{sk}}\) be a Skolem expansion of \(T^*\) as in
\cref{def:p03-morley-skolem-hull}, let \(\Phi\) be obtained from an
indiscernible sequence in a model of \(T^{*\mathrm{sk}}\) as in
\cref{thm:p03-w03-p03-028-coherent}, and for linear orders \(I\) and \(J\)
let \(H_I\) and \(H_J\) be term hulls of sequences \((b_i)_{i\in I}\) and
\((b'_j)_{j\in J}\) realizing \(\Phi\), as in
\cref{thm:p03-w03-p03-028-realize}.  For every order embedding
\(e\colon I\to J\) the assignment
\[
  t(b_{i_1},\ldots,b_{i_n})\longmapsto t(b'_{e(i_1)},\ldots,b'_{e(i_n)})
\]
is a well-defined embedding of \(H_I\) into \(H_J\) in the Skolem language,
and its restriction to the \(\Lang\)-reducts is an elementary embedding.
If \(e\) is onto, it is an isomorphism.  In particular any two hulls built
over the same \(J\) are isomorphic.
\end{hblemma}

\section{Ehrenfeucht--Mostowski constructions and embedding properties}
\label{sec:p03-w03-p03-028-em}

Ehrenfeucht and Mostowski stated their main results in terms of
automorphism groups.  For a set \(X\) linearly ordered by \(\prec\), let
\(G(X,\prec)\) be the group of all transformations of \(X\) onto itself
that preserve \(\prec\) \cite[p.~63]{EhrenfeuchtMostowski1956}.  A group
\(G_1\) of transformations of \(X_1\) \emph{strongly contains} a group
\(G\) of transformations of \(X\) when \(X\subseteq X_1\) and every
\(f\in G\) extends to some \(f_1\in G_1\)
\cite[p.~59]{EhrenfeuchtMostowski1956}.  For a model \(\mathcal M\),
\(G_{\mathcal M}\) is its automorphism group
\cite[p.~56]{EhrenfeuchtMostowski1956}.

\begin{hbtheorem}{thm:p03-w03-p03-028-em-57}{proved=annex:p03-w03-p03-028-em-57}{Ehrenfeucht--Mostowski Theorem 5.7}
If a theory \(S\) has at least one model over an infinite set, then for
every ordered set \((X,\prec)\) there is a model \(\mathcal M_0\) of \(S\)
such that \(G_{\mathcal M_0}\) strongly contains \(G(X,\prec)\)
\cite[Theorem~5.7, p.~63]{EhrenfeuchtMostowski1956}.
\end{hbtheorem}

\begin{hbcorollary}{thm:p03-w03-p03-028-em-56}{proved=annex:p03-w03-p03-028-em-56}{Ehrenfeucht--Mostowski Theorem 5.6}
Every theory \(S\) that has at least one model over an infinite set has a
model \(\mathcal M_0\) such that \(G_{\mathcal M_0}\) strongly contains an
infinite cyclic group \cite[Theorem~5.6, p.~60]{EhrenfeuchtMostowski1956}.
\end{hbcorollary}

Ehrenfeucht and Mostowski give an independent proof of Theorem~5.6, because
it is much simpler, while noting that it is a special case of Theorem~5.7
\cite[p.~60]{EhrenfeuchtMostowski1956}.  Both of their proofs use the
partition theorem once, in the consistency argument: for Theorem~5.6 on the
\((2n+1)\)-element subsets of an infinite model with \(2^s\) classes
\cite[p.~62]{EhrenfeuchtMostowski1956}, and for Theorem~5.7 on the
\(n\)-element subsets with \(2^s\) classes, where \(n\) is the number of
generators in the finite fragment and \(s\) the number of matrices
\cite[pp.~64--65]{EhrenfeuchtMostowski1956}.  The annex proves both
statements twice: through \cref{thm:p03-morley-ramsey-em}, and by the
article's own argument, which is reconstructed in
\cref{sec:p03-w03-p03-028-em1956,sec:p03-w03-p03-028-em-automorphisms,sec:p03-w03-p03-028-em-induced,sec:p03-w03-p03-028-em-original}.

\begin{hbexample}{def:p03-w03-p03-028-shift-example}{The shift on the integers}
Take \(X=\mathbb Z\) with its usual order and \(h(x)=x+1\).  The powers of
\(h\) form an infinite cyclic subgroup of \(G(\mathbb Z,<)\), since
\(h^n(0)=n\neq0\) for \(n\neq0\).  In a model generated by
order-indiscernible copies of \(\mathbb Z\), the automorphism induced by
\(h\) moves every generator one step up, and its powers are distinct
because they already differ on the generators.  This is the special case
used for Theorem~5.6.
\end{hbexample}

\begin{hbtheorem}{thm:p03-w03-p03-028-no-finite-order}{proved=annex:p03-w03-p03-028-no-finite-order}{Infinite cyclic cannot be replaced by an arbitrary group}
Let the theory \(S\) contain a binary predicate \(\pi\) with the axioms
that \(\pi\) is transitive and irreflexive and that any two elements
\(x,y\) satisfy \(x\mathrel\pi y\), \(x=y\) or \(y\mathrel\pi x\).  Let
\(\mathcal M\) be a model of \(S\) over a set \(X_1\), let
\(X\subseteq X_1\), and let \(f\) be an automorphism of \(\mathcal M\)
that maps \(X\) onto itself.  If the restriction of \(f\) to \(X\) has
finite order, then that restriction is the identity.  Hence, for every
model \(\mathcal M\) of \(S\), \(G_{\mathcal M}\) strongly contains no
group of transformations of a subset of its universe that has an element of
finite order other than the identity
\cite[pp.~62--63]{EhrenfeuchtMostowski1956}.
\end{hbtheorem}

In particular the conclusion ``an infinite cyclic group'' in
\cref{thm:p03-w03-p03-028-em-56} cannot be strengthened to ``an arbitrary
transformation group''.  The order-preserving group \(G(X,\prec)\) in
Theorem~5.7 has no non-identity element of finite order, so the two
statements are consistent.

\subsection{The 1956 article: theories, pseudo-models and the model \texorpdfstring{\(\mathfrak M(X,I)\)}{M(X,I)}}
\label{sec:p03-w03-p03-028-em1956}

The theorems above are stated in the form Ehrenfeucht and Mostowski gave
them, but proved in the annex along Morley's later route.  This subsection
and the next three reconstruct the article itself, in its own notation and
with its own proofs \cite[pp.~50--68]{EhrenfeuchtMostowski1956}.  The article
studies models of theories in first-order logic with identity and their
automorphisms.  Its main results are in its section~5, chiefly Theorem~5.7,
and one corollary is that axiom systems of arithmetic have models with
non-trivial automorphisms, which answers a question of G.~Hasenjaeger
\cite[p.~50]{EhrenfeuchtMostowski1956}.  The proofs are non-constructive in
two places: a consistent theory is extended to a consistent and complete set
of matrices, and every set is assumed to admit a linear ordering (the
ordering principle).  Because the theories may have any number of constants,
Ehrenfeucht and Mostowski observe that both tools reduce to the fundamental
theorem on ideals in Boolean algebras, and that the only proof using the full
axiom of choice is that of Theorem~4.3
\cite[pp.~50--51]{EhrenfeuchtMostowski1956}.  They describe the method as an
instance of what Tarski called the condensation of singularities: a model
with a large automorphism group exists exactly when infinitely many
sentences are satisfiable together, and this is shown by checking that any
finitely many of them can be added to the axioms without inconsistency
\cite[p.~50]{EhrenfeuchtMostowski1956}.  Sections~1--3 of the article
collect notions and lemmas and contain no new results; the construction of
models in section~3 follows Henkin, Novak and Rasiowa
\cite[p.~51]{EhrenfeuchtMostowski1956}.

\begin{hbconvention}{def:p03-w03-p03-028-em-syntax}{The syntax of the 1956 article}
A theory \(S\) is given by three sets of arbitrary cardinality: the set
\(F(S)\) of functors, the set \(P(S)\) of predicates, and the set \(A(S)\) of
axioms \cite[p.~51]{EhrenfeuchtMostowski1956}.  \(P(S)\) contains the
identity predicate \(\iota\).  For a functor or predicate \(\varphi\),
\(a(\varphi)\) is its number of arguments; \(a(\iota)=2\), and a functor with
\(a(\varphi)=0\) is a \emph{constant}.  All theories share the variables
\(\xi_1,\xi_2,\ldots\).  \(W(S)\), the class of \emph{terms}, is the least
class containing the variables and containing
\(\varphi(\omega_1,\ldots,\omega_{a(\varphi)})\) for \(\varphi\in F(S)\)
whenever it contains \(\omega_1,\ldots,\omega_{a(\varphi)}\).  \(Z(S)\), the
class of \emph{matrices}, is the least class containing
\(\pi(\omega_1,\ldots,\omega_{a(\pi)})\) for \(\pi\in P(S)\) and terms
\(\omega_j\), and closed under \(\mathord\sim\zeta\), \(\zeta_1\cdot\zeta_2\)
and \((\exists\xi_n)\zeta\); the other connectives (\(\supset\), \(\equiv\),
\(\vee\)) are defined from these \cite[p.~51]{EhrenfeuchtMostowski1956}.
\(\operatorname{subst}\alpha\binom{\xi_1,\xi_2,\ldots}{\omega_1,\omega_2,\ldots}\),
or \(\operatorname{subst}\alpha(\omega_1,\omega_2,\ldots)\), is the result of
substituting \(\omega_j\) for \(\xi_j\) in a term or matrix \(\alpha\), where
the substitution is performable \cite[p.~52]{EhrenfeuchtMostowski1956}.  A
matrix is \emph{open} if it has no bound variables and \emph{closed} if it has
no free variables; a term without variables is \emph{constant}, and \(W^*(S)\)
is the set of constant terms.  \(T(S)\) is the class of theorems of \(S\); the
matrices
\(\xi_k\iota\xi_k\), \(\xi_k\iota\xi_l\supset\xi_l\iota\xi_k\),
\((\xi_k\iota\xi_l)\cdot(\xi_l\iota\xi_m)\supset(\xi_k\iota\xi_m)\),
\(\xi_k\iota\xi_l\supset\omega\iota\operatorname{subst}\omega\binom{\xi_k}{\xi_l}\)
and
\(\xi_k\iota\xi_l\supset[\zeta\equiv\operatorname{subst}\zeta\binom{\xi_k}{\xi_l}]\)
belong to \(T(S)\) for every \(S\), every term \(\omega\) and every matrix
\(\zeta\) \cite[p.~52]{EhrenfeuchtMostowski1956}.  A theory \(S'\) is an
\emph{extension} of \(S\) if \(F(S)\subseteq F(S')\), \(P(S)\subseteq P(S')\)
and \(T(S)\subseteq T(S')\); the extension is \emph{inessential} if
\(P(S')=P(S)\) and \(T(S')\cap Z(S)=T(S)\).  \(S\) is \emph{open} if every
axiom is an open matrix \cite[p.~52]{EhrenfeuchtMostowski1956}.

The article does not fix a deductive calculus.  Here \(S\) is read as a
first-order signature with \(\iota\) as the equality symbol, \(\mathord\sim\),
\(\cdot\) and \(\exists\) as \(\neg\), \(\land\) and \(\exists\), and
\(T(S)\) as the set of matrices whose universal closure is derivable in
\(\mathsf{ND}_{\mathsf{FOL}}\) from the universal closures of the members of
\(A(S)\).  \(S\) is \emph{consistent} when these closures form a consistent
set (\cref{sec:fol-metatheory-01}).  With this reading the five
identity matrices above are theorems of every theory, and
\cref{thm:p03-w03-p03-028-em-theorems-valid} shows that \(T(S)\) is the same
for every sound and complete calculus.
\end{hbconvention}

\begin{hbdefinition}{def:p03-w03-p03-028-em-pseudo-model}{Pseudo-models, satisfaction and models}
Let \(S\) be a theory and \(X\) a non-empty set.  A \emph{pseudo-model} \(M\) of
\(S\) over \(X\) assigns to each \(\varphi\in F(S)\) a function \(M_\varphi\)
of \(a(\varphi)\) arguments from \(X\) to \(X\), to each \(\pi\in P(S)\) an
\(a(\pi)\)-ary relation \(M_\pi\) on \(X\), and to \(\iota\) the identity of
\(X\) \cite[p.~52]{EhrenfeuchtMostowski1956}; in the handbook's terms it is
an \(S\)-structure with domain \(X\).  A \emph{valuation} is a function \(f\)
with \(f(\xi_j)=x_j\in X\).  The value \(\operatorname{val}_{fM}\omega\) of a
term is defined by \(\operatorname{val}_{fM}\xi_j=x_j\) and
\(\operatorname{val}_{fM}\varphi(\omega_1,\ldots,\omega_{a(\varphi)})
=M_\varphi(\operatorname{val}_{fM}\omega_1,\ldots,
\operatorname{val}_{fM}\omega_{a(\varphi)})\), and is also written
\(\operatorname{val}_M\omega(x_1,x_2,\ldots)\)
\cite[pp.~52--53]{EhrenfeuchtMostowski1956}.  Satisfaction
\(\operatorname{stsf}_{fM}\zeta\), also written
\(\operatorname{stsf}_M\zeta(x_1,x_2,\ldots)\), holds for an atomic matrix
\(\pi(\omega_1,\ldots,\omega_{a(\pi)})\) iff the relation \(M_\pi\) holds of
the values \(\operatorname{val}_{fM}\omega_j\); it commutes with
\(\mathord\sim\) and \(\cdot\); and \(\operatorname{stsf}_{fM}(\exists\xi_n)\zeta\)
holds iff \(\operatorname{stsf}_{f'M}\zeta\) for some \(f'\sim_nf\), where
\(f'\sim_nf\) means \(f'(\xi_j)=f(\xi_j)\) for \(j\neq n\)
\cite[p.~53]{EhrenfeuchtMostowski1956}.  \(V_M\) is the set of matrices
\(\zeta\) with \(\operatorname{stsf}_{fM}\zeta\) for every valuation \(f\), and
\(M\) is a \emph{model} of \(S\) if \(A(S)\subseteq V_M\).  If \(S'\) is an
extension of \(S\) and \(M'\), \(M\) are pseudo-models of \(S'\), \(S\) over
the same \(X\), then \(M'\) is an \emph{extension} of \(M\) if
\(M'_\varphi=M_\varphi\) for every \(\varphi\in F(S)\cup P(S)\)
\cite[p.~53]{EhrenfeuchtMostowski1956}.
\end{hbdefinition}

The article allows models over arbitrary sets; the handbook's semantics uses
non-empty domains, and every model below has a non-empty domain.

\begin{hbtheorem}{thm:p03-w03-p03-028-em-21}{proved=annex:p03-w03-p03-028-em-21}{Ehrenfeucht--Mostowski Theorem 2.1}
If \(S'\) is an extension of \(S\) and \(M'\) an extension of \(M\), then
\(\operatorname{stsf}_{fM'}\zeta\equiv\operatorname{stsf}_{fM}\zeta\) for
every valuation \(f\) and every \(\zeta\in Z(S)\)
\cite[Theorem~2.1, p.~53]{EhrenfeuchtMostowski1956}.
\end{hbtheorem}

\begin{hblemma}{thm:p03-w03-p03-028-em-theorems-valid}{proved=annex:p03-w03-p03-028-em-theorems-valid}{Theorems are the matrices valid in all models}
Assume the axiom of choice.  For a theory \(S\) and \(\zeta\in Z(S)\),
\(\zeta\in T(S)\) iff \(\zeta\in V_M\) for every model \(M\) of \(S\).  \(S\)
is consistent iff it has a model.
\end{hblemma}

\begin{hbtheorem}{thm:p03-w03-p03-028-em-11}{proved=annex:p03-w03-p03-028-em-11}{Ehrenfeucht--Mostowski Theorem 1.1}
Assume the axiom of choice.  For every theory \(S\) there is an open theory
\(S'\) which is an inessential extension of \(S\)
\cite[Theorem~1.1, p.~52]{EhrenfeuchtMostowski1956}.
\end{hbtheorem}

Ehrenfeucht and Mostowski obtain Theorem~1.1 from the second
\(\varepsilon\)-theorem of Hilbert and Bernays
\cite[p.~52]{EhrenfeuchtMostowski1956}.  The proof in the annex takes \(S'\) to
be a Skolemization (\cref{thm:skolem-equisatisfiable}) and uses the completeness
theorem, so it assumes the axiom of choice, a hypothesis the article's
statement does not carry; it also shows that every model of \(S\) expands to
a model of this \(S'\).

\begin{hbdefinition}{def:p03-w03-p03-028-em-mxi}{The theories \(S^*(X)\), \(S(X,I)\) and the pseudo-model \(\mathfrak M(X,I)\)}
Let \(S\) be an open theory with a model over an infinite set, and let \(X\)
be a set of symbols not occurring in \(S\), identified with the elements of a
given set \cite[p.~53]{EhrenfeuchtMostowski1956}.  \(S^*(X)\) is obtained by
adding the elements of \(X\) to \(F(S)\) as constants (\(a(x)=0\)) and the
matrices \(\mathord\sim(x'\iota x'')\), for \(x',x''\in X\) with
\(x'\neq x''\), to \(A(S)\) \cite[p.~53]{EhrenfeuchtMostowski1956}.  Let
\(I\) be a consistent and complete subset of \(Z(S^*(X))\) containing
\(A(S^*(X))\); here this means that the theory \(S(X,I)\) with
\(F(S(X,I))=F(S^*(X))\), \(Z(S(X,I))=Z(S^*(X))\) and \(A(S(X,I))=I\) is
consistent, that \(T(S(X,I))\subseteq I\), and that for every closed
\(\zeta\in Z(S^*(X))\) either \(\zeta\in I\) or \(\mathord\sim\zeta\in I\)
\cite[p.~54]{EhrenfeuchtMostowski1956}.  Constant terms \(\omega_1,\omega_2\)
of \(S(X,I)\) are \emph{equivalent}, \(\omega_1\approx\omega_2\), if
\(\omega_1\iota\omega_2\in I\).  \(\mathfrak X_X\) is the set of
\(\approx\)-classes of \(W^*(S(X,I))\), and \([\omega]\) is the class of
\(\omega\).  For \(\varphi\in F(S)\) and \(\pi\in P(S)\) put
\cite[p.~54]{EhrenfeuchtMostowski1956}
\begin{gather*}
  \mathfrak M_\varphi([\omega_1],\ldots,[\omega_{a(\varphi)}])
   =[\varphi(\omega_1,\ldots,\omega_{a(\varphi)})],\\
  \mathfrak M_\pi([\omega_1],\ldots,[\omega_{a(\pi)}])
   \equiv\pi(\omega_1,\ldots,\omega_{a(\pi)})\in I .
\end{gather*}
The resulting \(\mathfrak M\) is written \(\mathfrak M(X,I)\)
\cite[p.~55]{EhrenfeuchtMostowski1956}.  For \(\mathfrak X_X\) to be non-empty,
\(X\) must be non-empty or \(F(S)\) must contain a constant; this side
condition is implicit in the article and is assumed whenever
\(\mathfrak M(X,I)\) is formed.
\end{hbdefinition}

The reading of ``complete'' above, which includes \(T(S(X,I))\subseteq I\), is
the one the article's Lemmas~3.2--3.6 use: they treat membership in \(I\) as
closed under the identity theorems, substitution and the propositional rules.

\begin{hblemma}{thm:p03-w03-p03-028-em-31}{proved=annex:p03-w03-p03-028-em-31}{Ehrenfeucht--Mostowski Lemma 3.1}
The theory \(S^*(X)\) is consistent
\cite[Lemma~3.1, p.~54]{EhrenfeuchtMostowski1956}.
\end{hblemma}

\begin{hblemma}{thm:p03-w03-p03-028-em-complete-set}{proved=annex:p03-w03-p03-028-em-complete-set}{Complete sets exist}
Assume the axiom of choice.  If a theory \(S_1\) is consistent, there is a
set \(I\subseteq Z(S_1)\) with \(A(S_1)\subseteq I\) such that the theory with
the symbols of \(S_1\) and axioms \(I\) is consistent, has all its theorems in
\(I\), and has \(\zeta\in I\) or \(\mathord\sim\zeta\in I\) for each closed
\(\zeta\in Z(S_1)\).  In particular a consistent and complete \(I\) as in
\cref{def:p03-w03-p03-028-em-mxi} exists
\cite[p.~54]{EhrenfeuchtMostowski1956}.
\end{hblemma}

\begin{hblemma}{thm:p03-w03-p03-028-em-32}{proved=annex:p03-w03-p03-028-em-32}{Ehrenfeucht--Mostowski Lemmas 3.2 and 3.3}
The relation \(\approx\) is an equivalence relation, and
\(x_1\not\approx x_2\) for distinct \(x_1,x_2\in X\).  If
\(\varphi\in F(S(X,I))\), \(\pi\in P(S(X,I))\), and
\(\omega_j\approx\omega'_j\), \(\tau_k\approx\tau'_k\) are constant terms
(\(j\leq a(\varphi)\), \(k\leq a(\pi)\)), then
\(\varphi(\omega_1,\ldots,\omega_{a(\varphi)})\approx
\varphi(\omega'_1,\ldots,\omega'_{a(\varphi)})\) and
\(\pi(\tau_1,\ldots,\tau_{a(\pi)})\in I\equiv
\pi(\tau'_1,\ldots,\tau'_{a(\pi)})\in I\)
\cite[Lemmas~3.2--3.3, p.~54]{EhrenfeuchtMostowski1956}.
\end{hblemma}

\begin{hblemma}{thm:p03-w03-p03-028-em-34}{proved=annex:p03-w03-p03-028-em-34}{Ehrenfeucht--Mostowski Lemmas 3.4--3.6}
\(\mathfrak M\) is a pseudo-model of \(S\) over \(\mathfrak X_X\)
\cite[Lemma~3.4, p.~55]{EhrenfeuchtMostowski1956}.  If \(\omega\in W(S)\),
\(\zeta\) is an open matrix of \(S\), and \(\tau_1,\tau_2,\ldots\) are
constant terms of \(S(X,I)\), then
\cite[Lemmas~3.5--3.6, p.~55]{EhrenfeuchtMostowski1956}
\begin{gather*}
  \operatorname{val}_{\mathfrak M}\omega([\tau_1],[\tau_2],\ldots)
   =[\operatorname{subst}\omega(\tau_1,\tau_2,\ldots)],\\
  \operatorname{stsf}_{\mathfrak M}\zeta([\tau_1],[\tau_2],\ldots)
   \equiv\operatorname{subst}\zeta(\tau_1,\tau_2,\ldots)\in I .
\end{gather*}
\end{hblemma}

\begin{hbtheorem}{thm:p03-w03-p03-028-em-37}{proved=annex:p03-w03-p03-028-em-37}{Ehrenfeucht--Mostowski Theorem 3.7}
\(\mathfrak M(X,I)\) is a model of \(S\) over \(\mathfrak X_X\)
\cite[Theorem~3.7, p.~55]{EhrenfeuchtMostowski1956}.
\end{hbtheorem}

\begin{hbtheorem}{thm:p03-w03-p03-028-em-38}{proved=annex:p03-w03-p03-028-em-38}{Ehrenfeucht--Mostowski Theorem 3.8}
Let \(S\) be any theory, \(X\) any set, \(M\) a pseudo-model of \(S\) over a
set \(Y\), and \(M'\) an extension of \(M\) to a pseudo-model of \(S^*(X)\)
over \(Y\).  If \(\omega\in W(S)\), \(\zeta\in Z(S)\) is open and
\(x_1,x_2,\ldots\in X\), then
\(\omega'=\operatorname{subst}\omega(x_1,x_2,\ldots)\) is a constant term and
\(\zeta'=\operatorname{subst}\zeta(x_1,x_2,\ldots)\) a closed matrix of
\(S^*(X)\), and \cite[Theorem~3.8, p.~55]{EhrenfeuchtMostowski1956}
\[
  \operatorname{val}_{M'}\omega'
   =\operatorname{val}_M\omega(M'_{x_1},M'_{x_2},\ldots),\qquad
  \operatorname{stsf}_{M'}\zeta'
   \equiv\operatorname{stsf}_M\zeta(M'_{x_1},M'_{x_2},\ldots).
\]
\end{hbtheorem}

\begin{hbwarning}{def:p03-w03-p03-028-em-open-only}{The need for open theories}
By \cref{thm:p03-w03-p03-028-em-34}, satisfaction in the model
\(\mathfrak M(X,I)\) agrees with membership in \(I\) for open matrices.
For matrices with quantifiers it can fail.  Let \(F(S)=\varnothing\),
\(P(S)=\{\iota\}\), \(A(S)=\varnothing\) and \(X=\{x\}\).  The set
\(A(S^*(X))\cup\{(\exists\xi_1)\mathord\sim(\xi_1\iota x)\}\) has a model with
two elements, so by \cref{thm:p03-w03-p03-028-em-complete-set} it lies in a
consistent and complete \(I\).  The only constant term is \(x\), so
\(\mathfrak X_X=\{[x]\}\) has one element, and the closed matrix
\((\exists\xi_1)\mathord\sim(\xi_1\iota x)\) of \(I\) is false in
\(\mathfrak M(X,I)\).  No term names the witness.  This is why the article
works with open theories and reaches arbitrary ones through Theorem~1.1,
whose added functors supply witnesses
(\cref{thm:p03-w03-p03-028-em-open-reduction}).
\end{hbwarning}

\subsection{Automorphisms and the representation of groups}
\label{sec:p03-w03-p03-028-em-automorphisms}

Let \(M\) be a model of \(S\) over \(X\).  A one-one mapping \(f\) of \(X\)
onto itself is an \emph{automorphism} of \(M\) if
\(f(M_\varphi(x_1,\ldots,x_{a(\varphi)}))=M_\varphi(f(x_1),\ldots,f(x_{a(\varphi)}))\)
and
\(M_\pi(x_1,\ldots,x_{a(\pi)})\equiv M_\pi(f(x_1),\ldots,f(x_{a(\pi)}))\)
for all \(\varphi\in F(S)\), \(\pi\in P(S)\) and \(x_1,x_2,\ldots\in X\)
\cite[p.~55]{EhrenfeuchtMostowski1956}; \(G_M\) is the group of automorphisms
of \(M\) \cite[p.~56]{EhrenfeuchtMostowski1956}.

\begin{hblemma}{thm:p03-w03-p03-028-em-41}{proved=annex:p03-w03-p03-028-em-41}{Ehrenfeucht--Mostowski Lemmas 4.1 and 4.2}
If \(f\in G_M\), \(\omega\in W(S)\), \(\zeta\in Z(S)\) and
\(x_1,x_2,\ldots\in X\), then
\(f(\operatorname{val}_M\omega(x_1,x_2,\ldots))
=\operatorname{val}_M\omega(f(x_1),f(x_2),\ldots)\) and
\(\operatorname{stsf}_M\zeta(x_1,x_2,\ldots)\equiv
\operatorname{stsf}_M\zeta(f(x_1),f(x_2),\ldots)\).  If \(S'\) is an
extension of \(S\), \(M\) a model of \(S\) over \(X\) and \(M'\) a model of
\(S'\) which is an extension of \(M\), then \(G_{M'}\subseteq G_M\)
\cite[Lemmas~4.1--4.2, p.~56]{EhrenfeuchtMostowski1956}.
\end{hblemma}

\begin{hbtheorem}{thm:p03-w03-p03-028-em-43}{proved=annex:p03-w03-p03-028-em-43}{Ehrenfeucht--Mostowski Theorem 4.3}
Assume the axiom of choice.  For each group \(G\) there is a theory \(S\) and
a model \(M\) of \(S\) such that the groups \(G_M\) and \(G\) are isomorphic
\cite[Theorem~4.3, p.~56, proof pp.~56--57]{EhrenfeuchtMostowski1956}.
\end{hbtheorem}

The theory in the proof has no functors; its predicates are \(\iota\) and
one binary predicate \(\pi_f\) for each one-one map \(f\) of \(G\) onto itself
that is not a left translation \(l(g)=g_0g\).  For such an \(f\) a pair
\(g_{1f},g_{2f}\) with \(f(g_{1f})g_{1f}^{-1}\neq f(g_{2f})g_{2f}^{-1}\) is
chosen, and \(M_{\pi_f}(g',g'')\) holds iff \(g'=gg_{1f}\) and \(g''=gg_{2f}\)
for some \(g\in G\) \cite[(4.3.1), p.~56]{EhrenfeuchtMostowski1956}.  The
simultaneous choice of these pairs is the use of the full axiom of choice
that the introduction of the article mentions.

\begin{hbexample}{def:p03-w03-p03-028-em-43-example}{Theorem~4.3 for the group of order 3}
Let \(G=\{0,1,2\}\) under addition modulo~3, so a left translation is
\(g\mapsto g_0+g\).  Of the six permutations of \(G\), the three translations
are automorphisms of the model built in the proof; the other three are the
transpositions.  For the transposition \(f\) exchanging \(1\) and \(2\),
\(f(0)-0=0\) and \(f(1)-1=1\), so one may take \(g_{1f}=0\) and
\(g_{2f}=1\).  Then \(M_{\pi_f}(g',g'')\) holds iff \(g''=g'+1\): the
relation \(M_{\pi_f}\) is the successor relation of the cycle
\(0\to1\to2\to0\).  The transposition \(f\) maps the pair \((1,2)\), which is
in the relation, to \((2,1)\), which is not, so \(f\) is not an automorphism.
The same argument applies to each transposition with its own predicate, so
\(G_M\) consists of the three translations and is isomorphic to \(G\).
\end{hbexample}

\subsection{Induced automorphisms and the consistency criterion}
\label{sec:p03-w03-p03-028-em-induced}

\begin{hbcorollary}{thm:p03-w03-p03-028-em-51}{proved=annex:p03-w03-p03-028-em-51}{Ehrenfeucht--Mostowski Theorem 5.1}
Let \(Y\) be an infinite set and \(Y^n\) the set of subsets of \(Y\) with
exactly \(n\) elements (\(n\geq1\)).  If
\(Y^n=C_1\cup\cdots\cup C_k\) is a partition of \(Y^n\) into mutually disjoint
sets, then there are \(j\leq k\) and an infinite \(Y_1\subseteq Y\) with
\(Y_1^n\subseteq C_j\) \cite[Theorem~5.1, p.~57]{EhrenfeuchtMostowski1956}.
\end{hbcorollary}

The article attributes Theorem~5.1 to Ramsey, citing ``theorem A on p.~384''
of an entry that its reference list gives as volume~30 (1929), pp.~338--384,
of the Proceedings of the London Mathematical Society
\cite[pp.~57, 68]{EhrenfeuchtMostowski1956}.  The printed article occupies
pp.~264--286 of volume~30 (1930), with Theorem~A on p.~264
(\cref{thm:p03-w03-p03-028-ramsey-a}).  For the rest of this subsection \(S\)
is an open theory and \(\mathfrak M(X,I)\) the model of \(S\) of
\cref{def:p03-w03-p03-028-em-mxi} \cite[p.~57]{EhrenfeuchtMostowski1956}.

\begin{hblemma}{thm:p03-w03-p03-028-em-52}{proved=annex:p03-w03-p03-028-em-52}{Ehrenfeucht--Mostowski Lemmas 5.2 and 5.3}
A one-one mapping \(h\) of \(X\) onto itself determines at most one
automorphism \(f\) of \(\mathfrak M(X,I)\) with \(f([x])=[h(x)]\) for
\(x\in X\).  If \(h\) has such an \(f\), \(h\) is said to \emph{induce an
automorphism}, and \(f\) is written \(f_h\).  If \(h_1\neq h_2\) and
\(f_{h_1}\), \(f_{h_2}\) exist, then \(f_{h_1}\neq f_{h_2}\)
\cite[Lemmas~5.2--5.3, p.~57]{EhrenfeuchtMostowski1956}.
\end{hblemma}

\begin{hblemma}{thm:p03-w03-p03-028-em-54}{proved=annex:p03-w03-p03-028-em-54}{Ehrenfeucht--Mostowski Lemma 5.4}
A one-one mapping \(h\) of \(X\) onto itself induces an automorphism of
\(\mathfrak M(X,I)\) iff for every open matrix \(\zeta\) of \(S\) and every
assignment \(\binom{\xi_1,\xi_2,\ldots}{x_1,x_2,\ldots}\) with
\(x_j\in X\) \cite[Lemma~5.4, p.~57, proof pp.~57--59]{EhrenfeuchtMostowski1956}
\begin{equation}
  \operatorname{subst}\zeta\binom{\xi_1,\xi_2,\ldots}{x_1,x_2,\ldots}
  \equiv
  \operatorname{subst}\zeta\binom{\xi_1,\xi_2,\ldots}{h(x_1),h(x_2),\ldots}
  \in I .
  \tag{5.4.1}
\end{equation}
In that case \(f_h([\tau])=[\tau^h]\), where \(\tau^h\) results from the
constant term \(\tau\) of \(S^*(X)\) by replacing each element \(x\) of \(X\)
occurring in it by \(h(x)\) \cite[p.~58]{EhrenfeuchtMostowski1956}.
\end{hblemma}

If \(G\) is the cyclic group generated by \(h\), the article says that
\(G_1\) strongly contains \(h\) when it strongly contains \(G\)
\cite[p.~60]{EhrenfeuchtMostowski1956}; strong containment is defined at the
start of \cref{sec:p03-w03-p03-028-em} \cite[p.~59]{EhrenfeuchtMostowski1956}.

\begin{hbtheorem}{thm:p03-w03-p03-028-em-55}{proved=annex:p03-w03-p03-028-em-55}{Ehrenfeucht--Mostowski Theorem 5.5}
Assume the axiom of choice.  Let \(S\) be an open theory and \(X\) a set.
There is a model \(M\) of \(S\) over a set \(X_1\supseteq X\) such that
\(G_M\) strongly contains a group \(G\) of transformations of \(X\) iff the
theory \(S^*(X)\) stays consistent when all equivalences
(5.4.1), with \(h\in G\) and \(\zeta\) an open matrix
of \(S\), are added to its axioms
\cite[Theorem~5.5, p.~60]{EhrenfeuchtMostowski1956}.
\end{hbtheorem}

\begin{hbexample}{def:p03-w03-p03-028-em-swap-example}{A permutation that induces no automorphism}
Let \(S\) have one binary predicate \(\pi\) besides \(\iota\) and the open
axioms \((\xi_1\pi\xi_2)\cdot(\xi_2\pi\xi_3)\supset(\xi_1\pi\xi_3)\),
\(\mathord\sim(\xi_1\pi\xi_1)\) and
\((\xi_1\pi\xi_2)\vee(\xi_1\iota\xi_2)\vee(\xi_2\pi\xi_1)\) of
\cref{thm:p03-w03-p03-028-no-finite-order}.  Let \(X=\{x_0,x_1\}\) and let
\(I\) contain \(x_0\pi x_1\); such an \(I\) exists by
\cref{thm:p03-w03-p03-028-em-complete-set,thm:p03-w03-p03-028-em-theorems-valid},
because a two-element order with \(x_0\) below \(x_1\) is a model.  For the
transposition \(h\) of \(x_0\) and \(x_1\), take \(\zeta=\xi_1\pi\xi_2\) and
an assignment with \(\xi_1\mapsto x_0\), \(\xi_2\mapsto x_1\).  Then
(5.4.1) reads \(x_0\pi x_1\equiv x_1\pi x_0\).  Its
left side is in \(I\) and its right side is not: together with
\(x_0\pi x_1\) and the transitivity axiom it would put \(x_0\pi x_0\) in
\(I\), against the instance \(\mathord\sim(x_0\pi x_0)\) of irreflexivity.  So
the equivalence is not in \(I\), and by \cref{thm:p03-w03-p03-028-em-54} \(h\)
induces no automorphism, in agreement with
\cref{thm:p03-w03-p03-028-no-finite-order}.
\end{hbexample}

\subsection{The original proofs of Theorems 5.6 and 5.7, and Theorem 6.1}
\label{sec:p03-w03-p03-028-em-original}

Both theorems are first proved for open theories, by showing that each
finite part of the set of equivalences in Theorem~5.5 is consistent with
\(S^*(X)\); a model of the finite part is built inside a given infinite
model, on a homogeneous set given by Theorem~5.1.

\begin{hblemma}{thm:p03-w03-p03-028-em-56-fragment}{proved=annex:p03-w03-p03-028-em-56-fragment}{The finite step in the proof of Theorem 5.6}
Let \(S\) be an open theory with a model over an infinite set, let
\(x_j\) (\(j\) an integer) be distinct new constants and \(h(x_j)=x_{j+1}\).
Let \(\zeta_1,\ldots,\zeta_s\) be open matrices of \(S\) with no variables
other than \(\xi_1,\ldots,\xi_t\), and let
\(i_l,j_l,\ldots,m_l\) (\(l=1,\ldots,s\)) be \(s\) sequences of \(t\)
integers in the interval \(-n\leq x\leq n\).  Let \(S^*\) have
\(P(S^*)=P(S)\), \(F(S^*)=F(S)\cup\{x_{-n},\ldots,x_{n+1}\}\), and as axioms
those of \(S\) together with
\cite[(5.6.1)--(5.6.2), p.~61]{EhrenfeuchtMostowski1956}
\begin{gather*}
  \operatorname{subst}\zeta_l\binom{\xi_1,\xi_2,\ldots,\xi_t}{x_{i_l},x_{j_l},\ldots,x_{m_l}}
  \equiv
  \operatorname{subst}\zeta_l\binom{\xi_1,\xi_2,\ldots,\xi_t}{h(x_{i_l}),h(x_{j_l}),\ldots,h(x_{m_l})}
  \quad(l\leq s),\\
  \mathord\sim(x_i\iota x_j)\quad(-n\leq i<j\leq n+1).
\end{gather*}
Then \(S^*\) has a model, and so is consistent
\cite[pp.~61--62]{EhrenfeuchtMostowski1956}.
\end{hblemma}

The proof orders the infinite model arbitrarily, divides its
\((2n+1)\)-element subsets into \(2^s\) classes by the sign pattern of
\(\zeta_1,\ldots,\zeta_s\) at the increasing enumeration, applies
Theorem~5.1, and interprets \(x_{-n},\ldots,x_{n+1}\) by \(2n+2\) increasing
elements of the homogeneous set; the two overlapping windows of \(2n+1\)
elements give the two sides of each equivalence.  The article speaks of
assigning \(2n+1\) different variables to the \(2n+2\) constants
\cite[p.~61]{EhrenfeuchtMostowski1956}; only the variables of
\(x_{-n},\ldots,x_n\) occur in its matrices (5.6.3), and the proof in the
annex uses exactly these.

\begin{hblemma}{thm:p03-w03-p03-028-em-57-fragment}{proved=annex:p03-w03-p03-028-em-57-fragment}{The finite step in the proof of Theorem 5.7}
Let \(S\) be an open theory with a model over an infinite set and let \(X\)
be ordered by \(\prec\).  Let \(\zeta_1,\ldots,\zeta_s\) be open matrices of
\(S\) with no variables other than \(\xi_1,\ldots,\xi_t\) (\(s,t\geq1\)), let
\(x_{(i-1)t+1},\ldots,x_{it}\) (\(i=1,\ldots,s\)) be elements of \(X\), not
necessarily distinct, let \(g_1,\ldots,g_s\in G(X,\prec)\), and let \(X^*\)
consist of these elements and all \(g_i(x_j)\) (\(i\leq s\), \(j\leq st\)).
Let \(S^*\) have \(F(S^*)=F(S)\cup X^*\), \(P(S^*)=P(S)\), and as axioms those
of \(S\) together with
\cite[(5.7.2)--(5.7.3), p.~63]{EhrenfeuchtMostowski1956}
\begin{gather*}
  \mathord\sim(x'\iota x'')\quad(x',x''\in X^*,\ x'\neq x''),\\
  \operatorname{subst}\zeta_i\binom{\xi_1,\ldots,\xi_t}{x_{(i-1)t+1},\ldots,x_{it}}
  \equiv
  \operatorname{subst}\zeta_i\binom{\xi_1,\ldots,\xi_t}{g_i(x_{(i-1)t+1}),\ldots,g_i(x_{it})}
  \quad(i\leq s).
\end{gather*}
Then \(S^*\) has a model, and so is consistent
\cite[pp.~63--66]{EhrenfeuchtMostowski1956}.
\end{hblemma}

Here the article divides the \(n\)-element subsets of the infinite model,
where \(n\) is the number of distinct elements among
\(x_1,\ldots,x_{st}\), according to the matrices
\(\psi_i=\operatorname{subst}\zeta_i(\bar x_{(i-1)t+1},\ldots,\bar x_{it})\)
evaluated at the \emph{distinguished ordering} of the subset: the unique
sequence \((u_1,\ldots,u_{st})\) of its elements with \(u_h\ll u_j\) iff
\(x_h\prec x_j\) \cite[pp.~64--65]{EhrenfeuchtMostowski1956}.  The images
\(g_i(x_1),\ldots,g_i(x_{st})\) are ordered like \(x_1,\ldots,x_{st}\), so the
homogeneous set gives both sides of each equivalence the same truth value
\cite[(5.7.8)--(5.7.10), pp.~65--66]{EhrenfeuchtMostowski1956}.

\begin{hblemma}{thm:p03-w03-p03-028-em-open-reduction}{proved=annex:p03-w03-p03-028-em-open-reduction}{Reduction to open theories}
Assume the axiom of choice.  Let \(S\) have a model over an infinite set and
let \(S'\) be an open inessential extension of \(S\) (Theorem~1.1).  Then
\(S'\) has a model over an infinite set.  If \(M'\) is a model of \(S'\) over
\(X_1\) and \(M\) is the pseudo-model of \(S\) over \(X_1\) of which \(M'\) is
an extension, then \(M\) is a model of \(S\) and \(G_{M'}\subseteq G_M\).
Hence, if \(G_{M'}\) strongly contains a group \(G\), so does \(G_M\), and
Theorems~5.6 and~5.7 for \(S\) follow from the same theorems for \(S'\)
\cite[pp.~62--63]{EhrenfeuchtMostowski1956}.
\end{hblemma}

The article states this step in one sentence: the general case reduces to the
open case by Theorems~1.1 and~4.2 \cite[p.~62]{EhrenfeuchtMostowski1956}; the
proof of Theorem~5.7 then restricts itself to open theories in the same way
\cite[p.~63]{EhrenfeuchtMostowski1956}.  The lemma supplies the fact the
sentence leaves implicit, that \(S'\) again has an infinite model.  With these
lemmas, Theorems~5.6 and~5.7
(\cref{thm:p03-w03-p03-028-em-56,thm:p03-w03-p03-028-em-57}) receive in the
annex the article's own proofs in addition to the proofs through
\cref{thm:p03-morley-ramsey-em}.

The article closes with a theorem that it describes as not directly connected
with its subject but needed in a later paper, with a proof close to that of
Theorem~5.7 \cite[pp.~66--67]{EhrenfeuchtMostowski1956}.

\begin{hbtheorem}{thm:p03-w03-p03-028-em-61}{proved=annex:p03-w03-p03-028-em-61}{Ehrenfeucht--Mostowski Theorem 6.1}
Assume the axiom of choice.  Let \(X\) be a set ordered by a relation
\(\prec\), and let \(S\) be an open theory with a model \(M\) over an infinite
set \(Y\) ordered by a relation \(\ll\).  Let \(Y^*\subseteq Y\) be infinite
and let \(\eta\) be an open matrix of \(S\) with the free variables
\(\xi_1,\ldots,\xi_q\) such that
\(\operatorname{stsf}_M\eta\binom{\xi_1,\ldots,\xi_q}{y_1,\ldots,y_q}\) for every
sequence \(y_1\ll y_2\ll\cdots\ll y_q\) of elements of \(Y^*\).  Then there is a
model \(M_0\) of \(S\) over a set \(X_1\supseteq X\) such that
\cite[Theorem~6.1, p.~67, proof pp.~67--68]{EhrenfeuchtMostowski1956}
\begin{enumerate}
\item[(6.1.1)] \(G_{M_0}\) strongly contains \(G(X,\prec)\), and
\item[(6.1.2)] \(\operatorname{stsf}_{M_0}\eta\binom{\xi_1,\ldots,\xi_q}{x_1,\ldots,x_q}\)
holds for every sequence \(x_1\prec x_2\prec\cdots\prec x_q\) in \(X\).
\end{enumerate}
\end{hbtheorem}

Theorem~6.1 adds to Theorem~5.7 the transfer of one open property: what an
open matrix says of all increasing \(q\)-tuples of an infinite set in the
given model, it says in \(M_0\) of all increasing \(q\)-tuples of \(X\).  It
does not make \(X\) a set of order indiscernibles; for that, every open
matrix true or false of increasing tuples has to be transferred at once, as
in \cref{thm:p03-w03-p03-028-finite-sat}.

\section{Nonisomorphic models of unstable theories}
\label{sec:p03-w03-p03-028-spectrum}

\begin{hbmethod}{def:p03-w03-p03-028-route}{The compactness, Ramsey and Skolem-hull route and its sources}
The route from a theory with an infinite model to Ehrenfeucht--Mostowski
models has the following steps.  For each step the list gives where the
handbook proves it and what each source contributes.
\begin{enumerate}
\item \emph{Partition theorem.}  Ramsey's Theorem~A,
\cref{thm:p03-w03-p03-028-ramsey-a}, from \cite[pp.~264--266]{Ramsey1930}.
Ehrenfeucht and Mostowski state it as their Theorem~5.1 for a finite
partition of the \(n\)-element subsets of an infinite set
\cite[Theorem~5.1, p.~57]{EhrenfeuchtMostowski1956}; they cite it as
Ramsey's Theorem~A ``on p.~384'' of an entry listed as pp.~338--384
\cite[p.~68]{EhrenfeuchtMostowski1956}, which does not match the page range
264--286 of the printed article, where Theorem~A is on p.~264
(\cref{thm:p03-w03-p03-028-em-51}).  Morley
states it as his Theorem~3.1 \cite[Theorem~3.1, p.~525]{Morley1965}.  The
finite Theorem~B, \cref{thm:p03-w03-p03-028-ramsey-b}, can replace it
(\cref{def:p03-w03-p03-028-finite-route}).
\item \emph{Partition step and finite satisfiability.}
\Cref{thm:p03-w03-p03-028-ramsey-step,thm:p03-w03-p03-028-finite-sat}: one
application of Theorem~A with arity the maximal tuple length \(N\) and
\(2^q\) classes.  Morley proves the same step with an arbitrary order of
the model, \(2^m\) classes and the same reduction to one length, by
embedding the index order in one without a largest element
\cite[Theorem~3.2, pp.~525--526]{Morley1965}.  In Ehrenfeucht and Mostowski
the step is the consistency of the extended theory \(S^*\) in the proofs of
Theorems~5.6 and~5.7 \cite[pp.~60--65]{EhrenfeuchtMostowski1956}.
\item \emph{Compactness.}  \Cref{thm:fol-compactness} turns finite
satisfiability into a model containing the indiscernibles
(\cref{thm:p03-ehrenfeucht-mostowski}).
\item \emph{Skolemization.}  Existential quantifiers are replaced by
function symbols, giving a universal conservative expansion; the proof is
in \cref{thm:p03-morley-ramsey-em}, following
\cite[Theorem~3.3, p.~526]{Morley1965}.  Ehrenfeucht and Mostowski work
with open theories and reduce the general case to that one by their
Theorems~1.1 and~4.2 \cite[p.~62]{EhrenfeuchtMostowski1956}; see
\cref{thm:p03-w03-p03-028-em-11,thm:p03-w03-p03-028-em-41,thm:p03-w03-p03-028-em-open-reduction}.
\item \emph{Term hull and induced maps.}  An order-preserving map of the
generators induces a well-defined embedding of the hull, and a bijection an
automorphism: \cref{thm:p03-morley-ramsey-em}, following
\cite[Theorems~3.4--3.5, p.~527]{Morley1965}; in the original, the
correspondence between generator maps and automorphisms of the model
\(\mathfrak M(X,I)\) of \cref{def:p03-w03-p03-028-em-mxi} is the content
of Lemmas~5.2--5.4 and Theorem~5.5
\cite[pp.~57--60]{EhrenfeuchtMostowski1956}, proved here as
\cref{thm:p03-w03-p03-028-em-52,thm:p03-w03-p03-028-em-54,thm:p03-w03-p03-028-em-55}.
\item \emph{Elementarity.}  With the theory first Morleyized, the induced
maps are elementary on the original reduct: \cref{thm:p03-morley-ramsey-em}
for maps inside one hull and \cref{thm:p03-w03-p03-028-elementary} for
maps between hulls over different orders,
following \cite[Theorem~3.6, p.~527]{Morley1965}.
\item \emph{Blueprints.}  Coherence of the blueprint and its realization
over any order:
\cref{thm:p03-w03-p03-028-coherent,thm:p03-w03-p03-028-realize}.
\end{enumerate}
The route ends with Ehrenfeucht--Mostowski models and the maps between
them.  The maximal-spectrum statement \cref{thm:p03-unstable-spectrum} is
not a consequence of these steps: it needs, in addition, a combinatorial
analysis of order-property formulas and a family of \(2^\kappa\) index
orders whose models can be told apart.  That theorem is cited to Shelah and
is not proved by the route above; \cref{sec:p03-w03-p03-028-shelah} and
\cref{def:p03-w03-p03-028-spectrum-boundary} record which of its steps
come from where, and \cref{thm:p03-w03-p03-028-two-models} gives the
weaker count that the route does prove.
\end{hbmethod}

\subsection{Two nonisomorphic models from the order property}
\label{sec:p03-w03-p03-028-two-models}

The route does give a nonisomorphic-model theorem, with a smaller count.
It compares two models of the same uncountable cardinality: an
Ehrenfeucht--Mostowski model over a well-order, in which few types are
realized over each countable set, and a model built by compactness around
an ordered family indexed by a dense order, in which many are.  The
hypothesis is the order property of \cref{def:p03-stability}; by that
definition a complete theory is unstable exactly when some formula has the
order property.

\emph{Inputs.}  A complete theory \(T\) in a countable language \(\Lang\);
the route of \cref{def:p03-w03-p03-028-route}, through
\cref{thm:p03-w03-p03-028-coherent,thm:p03-w03-p03-028-realize,thm:p03-w03-p03-028-elementary};
compactness (\cref{thm:fol-compactness}) and the downward
L\"owenheim--Skolem theorem (\cref{thm:downward-ls}), both under the axiom
of choice.  \emph{Conclusion.}  For every uncountable \(\kappa\), at least
two nonisomorphic models of \(T\) of cardinality \(\kappa\).

For a model \(\mathcal M\), a set \(Y\subseteq\dom\mathcal M\) and
\(n\geq1\), say that an \(n\)-type over \(Y\) is \emph{realized} in
\(\mathcal M\) when it is the set of \(\Lang\)-formulas
\(\psi(\bar x,\bar y)\) with parameters \(\bar y\) from \(Y\) satisfied by
some \(n\)-tuple \(\bar c\) of \(\mathcal M\); this is
\(\operatorname{tp}(\bar c/Y)\).

\begin{hblemma}{thm:p03-w03-p03-028-thin-tuples}{proved=annex:p03-w03-p03-028-thin-tuples}{Ehrenfeucht--Mostowski models over well-orders realize few types}
Let \(T\) be a complete theory in a countable language \(\Lang\) with an
infinite model, let \(T^{*\mathrm{sk}}\), \(\Phi\) and the hulls be as in
\cref{thm:p03-w03-p03-028-elementary}, and let \(\kappa\) be an infinite
cardinal, ordered as an ordinal.  Let \(H_\kappa\) be the term hull of a
sequence \((b_j)_{j<\kappa}\) realizing \(\Phi\), and \(\mathcal M\) its
\(\Lang\)-reduct.  Then \(\mathcal M\entails T\),
\(\abs{\dom\mathcal M}=\kappa\), and for every countable
\(Y\subseteq\dom\mathcal M\) and every \(n\geq1\) only countably many
\(n\)-types over \(Y\) are realized in \(\mathcal M\).
\end{hblemma}

For \(n=1\) this is the first clause of \cref{thm:p03-morley-thin-models};
the proof in the annex is written for tuples, which the next lemma needs,
and it makes explicit the one place where the well-ordering of the index
set is used: every element of \(\kappa\) has a least element of a given
countable set of indices above it, unless it lies above all of them.

\begin{hblemma}{thm:p03-w03-p03-028-thick}{proved=annex:p03-w03-p03-028-thick}{A model with many types over a countable set}
Let \(T\) be a complete theory in a countable language, and let
\(\varphi(\bar x;\bar y)\) have the order property in \(T\), with
\(\bar x\) of length \(n\).  For every uncountable cardinal \(\kappa\) there
are \(\mathcal N\entails T\) with \(\abs{\dom\mathcal N}=\kappa\) and a
countable \(Y\subseteq\dom\mathcal N\) over which at least \(\aleph_1\)
distinct \(n\)-types are realized in \(\mathcal N\).
\end{hblemma}

\begin{hbtheorem}{thm:p03-w03-p03-028-two-models}{proved=annex:p03-w03-p03-028-two-models}{Unstable countable theories are not categorical in uncountable cardinals}
Let \(T\) be a complete theory in a countable language in which some
formula has the order property (\cref{def:p03-stability}).  Then for every
uncountable cardinal \(\kappa\), the model \(\mathcal M\) of
\cref{thm:p03-w03-p03-028-thin-tuples} and the model \(\mathcal N\) of
\cref{thm:p03-w03-p03-028-thick} are models of \(T\) of cardinality
\(\kappa\) that are not isomorphic.  In particular \(T\) is not
\(\kappa\)-categorical.
\end{hbtheorem}

The proof uses only the route and the invariance of type counts under
isomorphism.  It does not pass through total transcendence, as the second
clause of \cref{thm:p03-morley-thin-models} does, and it does not use any
of Shelah's steps listed below.

\begin{hbexample}{def:p03-w03-p03-028-dlo-example}{Dense linear order}
Let \(T\) be the theory of dense linear orders without endpoints and
\(\varphi(x;y)\) the formula \(x<y\); with \(a_i=b_i=i\) in
\((\mathbb Q,<)\) it has the order property.  For \(\kappa=\aleph_1\) the
construction of \cref{thm:p03-w03-p03-028-thick} can be carried out by
hand: choose a set \(R\) of \(\aleph_1\) irrational reals and let
\(D=\mathbb Q\cup R\) with the order of \(\mathbb R\).  Between any two
reals lies a rational, so \(D\) is dense and has no endpoints, and
\(D\entails T\).  For \(r<r'\) in \(R\) choose \(q\in\mathbb Q\) with
\(r<q<r'\); the formula \(x<q\) belongs to the type of \(r\) over
\(\mathbb Q\) and not to that of \(r'\).  So \(D\) realizes \(\aleph_1\)
one-types over the countable set \(\mathbb Q\), while the model of
\cref{thm:p03-w03-p03-028-thin-tuples} over \(\omega_1\) realizes only
countably many over each countable subset.  The two models of cardinality
\(\aleph_1\) are not isomorphic.
\end{hbexample}

\begin{hbwarning}{def:p03-w03-p03-028-two-models-limits}{What the two-model theorem does not give}
\begin{enumerate}
\item \emph{Countable models.}  The theorem says nothing at
\(\kappa=\aleph_0\), and the conclusion fails there: the theory of
\cref{def:p03-w03-p03-028-dlo-example} has the order property and, by
the back-and-forth argument of \cref{thm:p03-rn-dlo-uniqueness}, exactly
one countable model up to isomorphism.  The proof breaks at \cref{thm:p03-w03-p03-028-thick}, since a
countable model cannot realize \(\aleph_1\) types.
\item \emph{The count.}  The theorem gives two models, not \(2^\kappa\).
The count \(2^\kappa\) of \cref{thm:p03-unstable-spectrum} needs the
further steps of Shelah's proof recorded below.
\item \emph{The language.}  The countability of \(\Lang\) is used in
\cref{thm:p03-w03-p03-028-thin-tuples}: the Skolem language has countably
many terms, so countably many patterns describe all types over a countable
set.  For an uncountable language the same argument bounds the number of
types over a set of size \(\abs{\Lang}\) by \(\abs{\Lang}\), and the
comparison model would need more than \(\abs{\Lang}\) types over a set of
that size; that comparison is not made here.
\item \emph{The definition of instability.}  The hypothesis is the order
property, which is how \cref{def:p03-stability} defines instability.
Shelah defines instability by counting types and obtains the order property
from it by a lemma quoted from his earlier work
\cite[Lemma~1.1, p.~475]{Shelah1971NonIsomorphicModels}.  The equivalence
with the handbook's definition follows from the stability spectrum theorem
\cref{thm:p03-w03-p03-032-spectrum}, whose proof supplies the counting-to-order
bridge used here.
\end{enumerate}
\end{hbwarning}

\subsection{Shelah's count and its dependencies}
\label{sec:p03-w03-p03-028-shelah}

Shelah's theorem is stated for pseudo-elementary classes.  For a complete
theory \(T\) in \(\Lang\) and a theory \(T_1\supseteq T\) in a language
\(\Lang_1\supseteq\Lang\), let \(I(\lambda,T_1,T)\) be the largest number
of pairwise nonisomorphic \(\Lang\)-reducts of models of \(T_1\) that have
cardinality \(\lambda\).  His main theorem says that if \(T\) is unstable
and \(\lambda>\abs{T_1}+\aleph_0\), then \(I(\lambda,T_1,T)=2^\lambda\)
\cite[Theorem~0.1, pp.~473--474]{Shelah1971NonIsomorphicModels};
\cref{thm:p03-unstable-spectrum} is the case \(T_1=T\).  Instability is
defined there by counting types: for every infinite \(\lambda\) some model
of \(T\) has a set \(A\) of at most \(\lambda\) elements over which more than
\(\lambda\) complete types are realized
\cite[p.~474]{Shelah1971NonIsomorphicModels}.  The proof has four parts.
\begin{enumerate}
\item \emph{Order property.}  An unstable \(T\) has a formula
\(\varphi(\bar x,\bar y)\) and sequences \(\bar a_0,\bar a_1,\ldots\) in a
model with \(\varphi(\bar a_m,\bar a_l)\) iff \(m<l\)
\cite[Lemma~1.1, p.~475]{Shelah1971NonIsomorphicModels}.  This lemma is
quoted from Shelah's earlier work, where it is Theorem~2.13, and is not
proved in the paper; by a coding of tuples the formula is then taken to be
binary and asymmetric, written \(x<y\)
\cite[Assumption~2, p.~475]{Shelah1971NonIsomorphicModels}.
\item \emph{Ehrenfeucht--Mostowski models.}  After \(T_1\) is made complete
and given Skolem functions, the construction of Ehrenfeucht and Mostowski
yields for every order \(I\) a model \(M_1(I)\) generated by an
indiscernible skeleton \((a_s)_{s\in I}\) on which \(a_s<a_t\) holds for
\(s<t\) \cite[Lemma~1.2 and Definition~1.1, p.~476]{Shelah1971NonIsomorphicModels}.
This is the only place where the route recorded in
\cref{def:p03-w03-p03-028-route} is used, and Shelah quotes it rather than
proving it.
\item \emph{Many contradictory orders.}  For every
\(\lambda>\aleph_0\) there are \(2^\lambda\) orders of cardinality
\(\lambda\), pairwise contradictory in the sense that no model contains
compatible nice sequences indexed by them
\cite[Definitions~2.1--2.2 and Theorem~2.2, pp.~477--481]{Shelah1971NonIsomorphicModels}.
This is pure order combinatorics, with no Ramsey argument.
\item \emph{Telling the models apart.}  A reduction to regular cardinals
\(\aleph_\alpha\) with \(\lambda>\aleph_\alpha>\abs{T_1}\)
\cite[Theorem~3.1 and Lemma~3.2, pp.~481--482]{Shelah1971NonIsomorphicModels},
and the proof that Ehrenfeucht--Mostowski models over orders built from
sets that differ modulo the closed-unbounded filter on \(\aleph_\alpha\) are
not isomorphic
\cite[Theorem~4.1 and Lemma~4.2, pp.~484--486]{Shelah1971NonIsomorphicModels}.
The last step uses that the skeleton of each such model is a nice sequence
and that an isomorphism carries nice sequences to nice sequences.
\end{enumerate}

\begin{proofboundary}{def:p03-w03-p03-028-spectrum-boundary}{What indiscernibles alone do not give}
Indiscernibles and Ehrenfeucht--Mostowski models give models with many
automorphisms (\cref{thm:p03-w03-p03-028-em-57}) and models that realize
few types over countable sets (\cref{thm:p03-morley-thin-models}).  They
do not by themselves give many nonisomorphic models: two different index
orders can yield isomorphic Ehrenfeucht--Mostowski models, for instance
when the theory is categorical in the relevant cardinal.  Together with the
order property and compactness they give two
(\cref{thm:p03-w03-p03-028-two-models}), and only for a countable language
and uncountable cardinals.  The count
\(2^\kappa\) in \cref{thm:p03-unstable-spectrum} rests on the order
property, on the combinatorics of contradictory orders and on the
nonisomorphism lemma listed above
\cite[Lemma~1.1, Theorem~2.2, Lemma~4.2]{Shelah1971NonIsomorphicModels};
the contradictory-order and nonisomorphism steps remain cited dependencies.
The counting-to-order bridge is proved in
\cref{thm:p03-w03-p03-032-spectrum}, and the construction of
Ehrenfeucht--Mostowski models is proved above.  Of the route in
\cref{def:p03-w03-p03-028-route}, only the existence of
Ehrenfeucht--Mostowski models with an ordered skeleton enters.  The
countable random graph, which is unstable and \(\aleph_0\)-categorical,
shows that the hypothesis \(\kappa>\abs T\) cannot be dropped.
\end{proofboundary}
